\documentclass{ctexart} 
% ctexbook  适合长文档：chapter、section、subsection、subsubsection
% ctexart   适合短文档：section、subsection、subsubsection

\usepackage{R3S-tex-doc}

% ---------- 自定义设置 ----------
% 目录分栏设置
% \multicoltoctrue      % 目录分两栏
% \multicoltocfalse     % 目录不分栏（默认）
\multicoltoctrue
% 颜色框自动编号设置
% \thmwithsectiontrue   % 颜色框随 section 编号（默认）
% \thmwithsectionfalse  % 颜色框全局编号
\thmwithsectionfalse
\colormodetrue
% 显示答案与解析设置
% \showanswerstrue      % 显示答案（默认）
% \showanswersfalse      % 不显示答案
% \showexplanationstrue   % 显示解析（默认）
% \showexplanationsfalse  % 不显示解析
\showanswerstrue
\showexplanationsfalse
% ----------------------------

% ctexbook放在前面会导致章节编号问题，ctexart放在前面则没有问题
\title{LaTeX 数学笔记模板} 
\author{RUINA}
\date{2026年8月1日}
% ---------- 正文 ----------
\begin{document}
% \frontmatter % 前置部分 (ctexart不支持)
\maketitle
\newpage
\tableofcontents % 目录
\newpage
% \mainmatter % 正文部分 (ctexart不支持)
\section{普通文本}
当第连的里大郭平屯杨略丑求清送事，句杀燕入本火土不陈云必大而不后司，亲召是疾反秦，
皇派真你斯无一不救杨侯反的同，土希把正尝不往而忧你作使马也的种，了长在但，光九非间洪松冒自君朗逃尘得谋头位，
入书尚办不的心回君到，于上葬极仅认不老苟作十下无说赏禀，不恨若韩。
\newline
Lorem ipsum dolor sit amet consectetur adipisicing elit. 
Autem deserunt eum cum ipsa inventore iure dignissimos praesentium ducimus facilis numquam optio, 
velit molestias. Quod nobis dolor quos itaque, cupiditate cumque!
\subsection{test}
\subsubsection{test1}

\begin{lstlisting}[language={[LaTeX]TeX}]
    当第连的里大郭平屯杨略丑求清送事，句杀燕入本火土不陈云必大而不后司，亲召是疾反秦，
皇派真你斯无一不救杨侯反的同，土希把正尝不往而忧你作使马也的种，了长在但，光九非间洪松冒自君朗逃尘得谋头位，
入书尚办不的心回君到，于上葬极仅认不老苟作十下无说赏禀，不恨若韩。
\newline
Lorem ipsum dolor sit amet consectetur adipisicing elit. 
Autem deserunt eum cum ipsa inventore iure dignissimos praesentium ducimus facilis numquam optio, 
velit molestias. Quod nobis dolor quos itaque, cupiditate cumque!
\end{lstlisting}



\newpage
\section{表格}
三线表
\begin{table}[htbp] % htbp用于定位
    \caption{表格演示}
    \centering %剧中
    \begin{tabular}{lccc}
        \toprule
        表头&列1&列2&列3 \\ % \\换行 &换列
        \midrule
        行1&text11&text12&text13 \\
        行2&text21&text22&text23 \\ 
        \bottomrule
    \end{tabular}        
\end{table}

\begin{center}
    \begin{tabular}{|c|c|c|}
        \hline
        1&2&3 \\ \hline
        4&5&6 \\ \hline
        7&8&9 \\ \hline
    \end{tabular}
\end{center}

\begin{center}
    \begin{tabular}{|c|c|c|}
        \hline
        1&2&3 \\ \hline
        \multicolumn{2}{|c|}{4} &5 \\ \hline
        6& \multicolumn{2}{|c|}{7} \\ \hline
        \multicolumn{3}{|c|}{8} \\ \hline
    \end{tabular}
\end{center}

\begin{center}
    \begin{tabular}{|l|c|c|c|}
        \toprule
        % 横向合并示例：表头跨3列
        \multicolumn{1}{c}{表头} & \multicolumn{3}{c}{合并的3列标题} \\
        \cmidrule(lr){2-4}  % 在第2-4列画横线
        表头 & 列1 & 列2 & 列3 \\
        \midrule
        
        % 纵向合并示例：行1-行2 的第一列合并
        \multirow{2}{*}{合并行} & text11 & text12 & text13 \\
                                & text21 & text22 & text23 \\
        
        \midrule
        
        % 横向+纵向同时合并的复杂示例
        \multirow{2}{*}{复杂合并} & \multicolumn{2}{c}{横向合并两列} & text13 \\
        \cmidrule(lr){2-3}
                                  & text21 & text22 & text23 \\
        
        \bottomrule
    \end{tabular}
\end{center}


\begin{lstlisting}[language={[LaTeX]TeX}]
    \begin{table}[htbp] % htbp用于定位
        \caption{表格演示}
        \centering %剧中
        \begin{tabular}{lccc}
            \toprule
            表头&列1&列2&列3 \\ % \\换行 &换列
            \midrule
            行1&text11&text12&text13 \\
            行2&text21&text22&text23 \\ 
            \bottomrule
        \end{tabular}        
    \end{table}
    
    \begin{center}
        \begin{tabular}{|c|c|c|}
            \hline
            1&2&3 \\ \hline
            4&5&6 \\ \hline
            7&8&9 \\ \hline
        \end{tabular}
    \end{center}
    
    \begin{center}
        \begin{tabular}{|c|c|c|}
            \hline
            1&2&3 \\ \hline
            \multicolumn{2}{|c|}{4} &5 \\ \hline
            6& \multicolumn{2}{|c|}{7} \\ \hline
            \multicolumn{3}{|c|}{8} \\ \hline
        \end{tabular}
    \end{center}

    \begin{tabular}{|l|c|c|c|}
        \toprule
        % 横向合并示例：表头跨3列
        \multicolumn{1}{c}{表头} & \multicolumn{3}{c}{合并的3列标题} \\
        \cmidrule(lr){2-4}  % 在第2-4列画横线
        表头 & 列1 & 列2 & 列3 \\
        \midrule
        
        % 纵向合并示例：行1-行2 的第一列合并
        \multirow{2}{*}{合并行} & text11 & text12 & text13 \\
                                & text21 & text22 & text23 \\
        
        \midrule
        
        % 横向+纵向同时合并的复杂示例
        \multirow{2}{*}{复杂合并} & \multicolumn{2}{c}{横向合并两列} & text13 \\
        \cmidrule(lr){2-3}
                                  & text21 & text22 & text23 \\
        
        \bottomrule
    \end{tabular}


\end{lstlisting}


\newpage
\section{图片}
\begin{center}
    % 测试图片过大所以调整大小
    \includegraphics[scale=0.2]{img/测试用JPG.jpg}
    \includegraphics[scale=0.2]{img/测试用PNG.png}      
\end{center}
\begin{lstlisting}[language={[LaTeX]TeX}]
    \begin{center}
        % 测试图片过大所以调整大小
        \includegraphics[scale=0.2]{img/测试用JPG.jpg}
        \includegraphics[scale=0.2]{img/测试用PNG.png}        
    \end{center}
\end{lstlisting}



\newpage
\section{列表}
    \begin{itemize}
        \item 战争
        \item 疾病
        \item 命运
        \item 死亡
    \end{itemize}

    当第连的里大郭平屯杨略丑求清送事，句杀燕入本火土不陈云必大而不后司，亲召是疾反秦，
皇派真你斯无一不救杨侯反的同，土希把正尝不往而忧你作使马也的种，了长在但，光九非间洪松冒自君朗逃尘得谋头位，
入书尚办不的心回君到，于上葬极仅认不老苟作十下无说赏禀，不恨若韩。

    \begin{itemize}
        \item test1
        \begin{itemize}
            \item 123
            \item 345
        \end{itemize}
        \item test2
    \end{itemize}
    
    文字分割

    \begin{enumerate}
        \item kaiser
        \item number
        \item lucy
    \end{enumerate}

\begin{lstlisting}[language={[LaTeX]TeX}]
    \begin{itemize}
        \item 战争
        \item 疾病
        \item 命运
        \item 死亡
    \end{itemize}

    当第连的里大郭平屯杨略丑求清送事，句杀燕入本火土不陈云必大而不后司，亲召是疾反秦，
皇派真你斯无一不救杨侯反的同，土希把正尝不往而忧你作使马也的种，了长在但，光九非间洪松冒自君朗逃尘得谋头位，
入书尚办不的心回君到，于上葬极仅认不老苟作十下无说赏禀，不恨若韩。

    \begin{itemize}
        \item test1
        \begin{itemize}
            \item 123
            \item 345
        \end{itemize}
        \item test2
    \end{itemize}
    
    文字分割

    \begin{enumerate}
        \item kaiser
        \item number
        \item lucy
    \end{enumerate}
\end{lstlisting}

\newpage
\section{代码}
% 旧版
\begin{verbatim}
test = {"lily":18,"lucy":19,"luck":20}

for key in test.keys():
    print(key)
    print(test[key])

for value in test.values():
    print(value)

for k,v in test.items():
    print(k,v)
\end{verbatim}

% 引入listings包
\begin{lstlisting}[language=python]
    test = {"lily":18,"lucy":19,"luck":20}

    for key in test.keys():
        print(key)
        print(test[key])

    for value in test.values():
        print(value)

    for k,v in test.items():
        print(k,v)
\end{lstlisting}

\newpage
\section{定理、定义等*}
%     定理类型1
%     \begin{theorem}{定理测试} 
%         这是一个测试文本
%     \end{theorem}

% \begin{lstlisting}[language={[LaTeX]TeX}]
%     定理类型1
%     \begin{theorem}{定理测试} 
%         这是一个测试文本
%     \end{theorem}
% \end{lstlisting}
    \begin{dfn}{测试}
    :这是一个测试定义
    \end{dfn}
    \begin{thm}{}
    :这是一个测试定理
    \end{thm}
    \begin{exm}{}
    :这是一个测试例子
    \end{exm}
    \begin{tip}{}   
    :这是一个测试提示
    \end{tip}
    \begin{warn}{}
    :这是一个测试错误
    \end{warn}
    \begin{exs}{}
    :这是一个测试练习
    \end{exs}
\begin{lstlisting}[language={[LaTeX]TeX}]
    % 在环境内容开始前加一个字符就正常显示，不然会吞掉一个字符
    \begin{dfn}{测试}
    :这是一个测试定义
    \end{dfn}
    \begin{thm}{}
    :这是一个测试定理
    \end{thm}
    \begin{exm}{}
    :这是一个测试例子
    \end{exm}
    \begin{tip}{}   
    :这是一个测试提示
    \end{tip}
    \begin{warn}{}
    :这是一个测试错误
    \end{warn}
    \begin{exs}{}
    :这是一个测试练习
    \end{exs}
\end{lstlisting}
\newpage
\section{公式*}
%行内
$ a_n \neq 0 $
%段落
\[
f(x)=a_n x^n + a_{n-1}x^{n-1}+\cdots + a_0
\]
\begin{lstlisting}[language={[LaTeX]TeX}]
    %行内
    $ a_n \neq 0 $
    %段落
    \[
    f(x)=a_n x^n + a_{n-1}x^{n-1}+\cdots + a_0
    \]
\end{lstlisting}

\begin{tabular}{|c|c|c|c|}
    \hline
    公式 & 代码 & 公式 & 代码 \\
    \hline
    $a_{b}$ & \verb|a_{b}| & $a^{b}$ & \verb|a^{b}| \\
    \hline
    $f'(x)$ & \verb|f'(x)| & $a_{b}^{c}$ & \verb|a_{b}^{c}| \\
    \hline
    $5/8,\frac{5}{8}$ & \verb|5/8,\frac{5}{8}| & $\sqrt[3]{2}$ & \verb|\sqrt[3]{2}| \\
    \hline   
\end{tabular}\\
以下为行内公式测试：\\
如图，线段 \(P_1 P_2\) 的端点 \(P_1\)，\(P_2\) 的坐标分别是$(x_1，y_1)$，$(x_2，y_2)$，
点 \(P\) 是直线 \(P_1 P_2\) 上异于 \(P_1\)，\(P_2\) 的点，当 \(\overrightarrow{P_1 P }=\lambda\overrightarrow{PP_2}\) 
($\lambda\neq 0$且 \(\lambda\neq 1\)) 时，\(\overrightarrow{OP} =\frac{1}{1 + \lambda} \overrightarrow{OP_1} + \frac{\lambda}{1 + \lambda} \overrightarrow{OP_2}\)，
点 \(P\) 的坐标是 \(\left(\frac{x_1 + \lambda x_2}{1 + \lambda} , \frac{y_1 + \lambda y_2}{1 + \lambda}\right)\)，
当 \(\lambda>0\) 时，点 \(P\) 在线段 \(P_1 P_2\) 上，称为内分点；当 \(\lambda<0\) 且 \(\lambda\neq -1\) 时，点 \(P\) 在线段 \(P_1 P_2\) 的延长线上，称为外分点。





\newpage
\section{练习与解答*}
演示：
\begin{exercise}
    求极限 $\displaystyle \lim_{x\to 0} \frac{\sin 3x}{x}$.
\end{exercise}
\begin{answer}
    3
\end{answer}
\begin{explanation}
    利用重要极限，$\displaystyle \lim_{x\to 0} \frac{\sin 3x}{x}=3\lim_{x\to 0}\frac{\sin 3x}{3x}=3\times 1=3$.
\end{explanation}

\begin{lstlisting}[language={[LaTeX]TeX}]
    \begin{exercise}
        求极限 $\displaystyle \lim_{x\to 0} \frac{\sin 3x}{x}$.
    \end{exercise}
    \begin{answer}
        3
    \end{answer}
    \begin{explanation}
        利用重要极限，$\displaystyle \lim_{x\to 0} \frac{\sin 3x}{x}=3\lim_{x\to 0}\frac{\sin 3x}{3x}=3\times 1=3$.
    \end{explanation}
\end{lstlisting}





\newpage
\section{分栏}
当第连的里大郭平屯杨略丑求清送事，句杀燕入本火土不陈云必大而不后司，亲召是疾反秦，
皇派真你斯无一不救杨侯反的同，土希把正尝不往而忧你作使马也的种，了长在但，光九非间洪松冒自君朗逃尘得谋头位，
入书尚办不的心回君到，于上葬极仅认不老苟作十下无说赏禀，不恨若韩。
\setlength{\columnseprule}{0.4pt}  % 栏间加竖线
\begin{multicols}{3}
    当第连的里大郭平屯杨略丑求清送事，句杀燕入本火土不陈云必大而不后司，亲召是疾反秦，
皇派真你斯无一不救杨侯反的同，土希把正尝不往而忧你作使马也的种，了长在但，光九非间洪松冒自君朗逃尘得谋头位，
入书尚办不的心回君到，于上葬极仅认不老苟作十下无说赏禀，不恨若韩。
\end{multicols}
当第连的里大郭平屯杨略丑求清送事，句杀燕入本火土不陈云必大而不后司，亲召是疾反秦，
皇派真你斯无一不救杨侯反的同，土希把正尝不往而忧你作使马也的种，了长在但，光九非间洪松冒自君朗逃尘得谋头位，
入书尚办不的心回君到，于上葬极仅认不老苟作十下无说赏禀，不恨若韩。


\begin{lstlisting}[language={[LaTeX]TeX}]
    当第连的里大郭平屯杨略丑求清送事，句杀燕入本火土不陈云必大而不后司，亲召是疾反秦，
皇派真你斯无一不救杨侯反的同，土希把正尝不往而忧你作使马也的种，了长在但，光九非间洪松冒自君朗逃尘得谋头位，
入书尚办不的心回君到，于上葬极仅认不老苟作十下无说赏禀，不恨若韩。
\setlength{\columnseprule}{0.4pt}  % 栏间加竖线
\begin{multicols}{3}
    当第连的里大郭平屯杨略丑求清送事，句杀燕入本火土不陈云必大而不后司，亲召是疾反秦，
皇派真你斯无一不救杨侯反的同，土希把正尝不往而忧你作使马也的种，了长在但，光九非间洪松冒自君朗逃尘得谋头位，
入书尚办不的心回君到，于上葬极仅认不老苟作十下无说赏禀，不恨若韩。
\end{multicols}
当第连的里大郭平屯杨略丑求清送事，句杀燕入本火土不陈云必大而不后司，亲召是疾反秦，
皇派真你斯无一不救杨侯反的同，土希把正尝不往而忧你作使马也的种，了长在但，光九非间洪松冒自君朗逃尘得谋头位，
入书尚办不的心回君到，于上葬极仅认不老苟作十下无说赏禀，不恨若韩。
\end{lstlisting}



\newpage
\section{绘图}
暂未学会，待补充
\section{测试1}\section{测试1}\section{测试1}\section{测试1}\section{测试1}
\section{测试1}\section{测试1}\section{测试1}\section{测试1}\section{测试1}
\section{测试1}\section{测试1}\section{测试1}\section{测试1}\section{测试1}
\section{测试1}\section{测试1}\section{测试1}\section{测试1}\section{测试10}
\subsection{测试用}
\noindent 段落不缩进
\par
段落缩进
\[
    E = m c^2
\]


\newpage
\appendix % 附录
\section{颜色}
\begin{table}[htbp]
    \centering
    \caption{颜色模式比较}
    \begin{tabular}{l|ccc}
        \toprule
        模式 & 语法 & 用途 & 取值范围 \\
        \midrule
        Named & \verb|\color{red}| & 简单快速 & 预定义名称 \\
        RGB & \verb|\color[RGB]{255,0,0}| & 屏幕显示 & 0-255 \\
        rgb & \verb|\color[rgb]{1,0,0}| & 屏幕显示 & 0-1 \\
        HTML & \verb|\color[HTML]{FF0000}| & Web/设计 & \#RRGGBB \\
        CMYK & \verb|\color[cmyk]{0,1,1,0}| & 印刷 & 0-1 \\
        HSB & \verb|\color[hsb]{0,1,1}| & 色彩调整 & 0-1 \\
        Gray & \verb|\color[gray]{0.5}| & 灰度 & 0-1 \\
        \bottomrule
    \end{tabular}
\end{table}
% \backmatter % 后置部分 (ctexart不支持)

\end{document}
% ---------------------------- 结束